\documentclass{resume}

\usepackage[left=0.75in,top=0.6in,right=0.75in,bottom=0.75in]{geometry}
\usepackage[utf8]{inputenc}
\usepackage[T1]{fontenc}

% font options
\usepackage[sc]{mathpazo}
\linespread{1.05}  % palladio needs more leading (space between lines)

\name{Maurício H. Vancine}
\address{State University of Campinas (Unicamp), Campinas, SP, Brazil}
\address{Institute of Biology, Department of Animal Biology, Rua Monteiro Lobato, 255, 13083-862}
\address{
\href{mailto:}{mauricio.vancine@gmail.com}
$\ \bullet\ $
\href{https://github.com/mauriciovancine}{GitHub}
$\ \bullet\ $
\href{https://www.linkedin.com/in/mauricio-vancine/}{LinkedIn}
$\ \bullet\ $
\href{https://mauriciovancine.github.io}{Website}}
\address{Ecologist, Data Scientist, and PhD in Ecology}

% cor dos links
\usepackage{xcolor}
\usepackage{hyperref}
\hypersetup{
    colorlinks=true,
    linkcolor=blue,
    urlcolor=blue,
    citecolor=blue
}

\begin{document}

%------------------------------------------------------------------------
% Career summary
%------------------------------------------------------------------------
\begin{rSection}{Career summary}
Environmental scientist and ecologist specializing in biodiversity conservation, spatial ecology, landscape ecology, and species distribution modeling. My research integrates ecological theory, geospatial analysis, and statistical and machine-learning approaches to investigate how landscape structure, habitat degradation, and climate change affect biodiversity patterns and ecosystem resilience. I also develop computational tools and datasets for ecological research and contribute to teaching and training in R.
\end{rSection}

%------------------------------------------------------------------------
% Academic expertise
%------------------------------------------------------------------------
\begin{rSection}{Academic expertise}
Spatial Ecology, Landscape Ecology, Species Distribution Modeling, Community and Functional Ecology, Geospatial Analysis (R and GIS), Ecological Statistics, Machine Learning for Ecology.
\end{rSection}

%------------------------------------------------------------------------
% Research interests
%------------------------------------------------------------------------
\begin{rSection}{Research interests}
Habitat loss, fragmentation and degradation; landscape metrics and connectivity; species distribution and joint species distribution models; amphibian ecology; biodiversity conservation; climate-change impacts; ecological data science and reproducible research.
\end{rSection}

%------------------------------------------------------------------------
% Education
%------------------------------------------------------------------------
\begin{rSection}{Education}
\textbf{São Paulo State University (UNESP)} \hfill {\em Mar 2020 - Jul 2024}\\
Ph.D. in Ecology, Evolution and Biodiversity\\
\textit{Dissertation}: Landscape structure as a predictor of taxonomic and functional diversity of amphibians in the Atlantic Forest\\
\textit{Advisor}: Prof. Dr. Milton Cezar Ribeiro; \href{https://hdl.handle.net/11449/256726}{\underline{dissertation}}

\textbf{São Paulo State University (UNESP)} \hfill {\em Mar 2016 - Jul 2018}\\
M.S. in Biological Sciences (Zoology)\\
\textit{Thesis}: Diversity, distribution and effect of climate change on Atlantic Forest amphibian communities\\
\textit{Advisor}: Prof. Dr. Célio Fernando Baptista Haddad; \href{http://hdl.handle.net/11449/154993}{\underline{thesis}}

\textbf{São Paulo State University (UNESP)} \hfill {\em Mar 2011 - Mar 2015}\\
B.S. in Ecology\\
\textit{Thesis}: Effect of fragmentation on the persistence of anuran amphibians within the Atlantic Forest\\
\textit{Advisor}: Prof. Dr. Milton Cezar Ribeiro; \href{http://hdl.handle.net/11449/138991}{\underline{thesis}}
\end{rSection}

%------------------------------------------------------------------------
% Professional experience
%------------------------------------------------------------------------
\begin{rSection}{Professional experience}
\begin{itemize}
\item \textbf{Postdoctoral researcher} \hfill{\em 2024-present}\\
\textit{Institution}: State University of Campinas (Unicamp), Landscape Diversity Laboratory (LDDiv)\\
\textit{Project}: Impacts of edge effects, fire, and climate change on vegetation composition and functional diversity across spatial scales in the Amazon rainforest.\\
\textit{Supervisor}: Prof. Dr. Mathias Mistretta Pires

\item \textbf{Postdoctoral research fellow / visiting researcher} \hfill{\em 2026}\\
\textit{Institution}: Yale School of the Environment, Brando Lab, USA\\
\textit{Project}: Effects of habitat degradation on inferring the local dominance of tree species in the Amazon rainforest.\\
\textit{Supervisor}: Prof. Dr. Paulo Brando

\item \textbf{Research assistant} \hfill{\em 2015-2016}\\
\textit{Institution}: São Paulo State University (UNESP), Spatial Ecology and Conservation Lab (LEEC)\\
\textit{Role}: Statistical and spatial analyses of mammal and ant biodiversity in the Atlantic Forest. Supervisor: Prof. Dr. Milton Cezar Ribeiro.

\item \textbf{Research intern} \hfill{\em 2014}\\
\textit{Institution}: University of São Paulo (USP), Ecology, Management and Conservation of Wild Fauna Lab (LEMaC)\\
\textit{Role}: Species distribution modeling. Supervisor: Prof. Dr. Katia Maria Paschoaletto Micchi de Barros Ferraz.
\end{itemize}
\end{rSection}

%------------------------------------------------------------------------
% Research experience
%------------------------------------------------------------------------
\begin{rSection}{Research experience}
\begin{itemize}
\item \textbf{Amazon rainforest (postdoctoral research, 2024-present):} investigates how edge effects, fire, and climate change influence vegetation composition, functional diversity, and the potential distribution of vertebrate fauna across spatial scales.
\item \textbf{Atlantic Forest (doctoral research, 2020-2024):} characterized landscape structure and evaluated its relationships with taxonomic and functional diversity of amphibian communities, including species-level patterns and co-occurrences using joint species distribution models.
\item \textbf{Atlantic Forest (master's research, 2016-2018):} assessed amphibian community surveys and investigated the effects of habitat loss, fragmentation, and climate change on amphibian persistence, including the genus \textit{Brachycephalus}.
\item \textbf{Atlantic Forest (undergraduate research, 2013-2015):} evaluated the contribution of landscape attributes, including forest cover, connectivity, relief, and urban proximity, to anuran persistence using species distribution modeling.
\end{itemize}
\end{rSection}

%------------------------------------------------------------------------
% Funding
%------------------------------------------------------------------------
\begin{rSection}{Funding}
\begin{itemize}
\item \textbf{Postdoctoral scholarship}, São Paulo Research Foundation (FAPESP), 2024-present. Grant 2024/19865-6. \href{https://bv.fapesp.br/en/bolsas/223203/impacts-of-edge-effects-fire-and-climate-change-on-vegetation-composition-and-functional-diversity-a/}{\underline{Project link}}.
\item \textbf{Doctoral scholarship}, FAPESP, 2022-2024. Grant 2022/01899-6. \href{https://bv.fapesp.br/en/bolsas/203713/landscape-structure-as-a-predictor-of-taxonomic-and-functional-diversity-of-amphibians-in-the-atlant/}{\underline{Project link}}.
\item \textbf{Master's scholarship}, FAPESP, 2017-2018. Grant 2017/09676-8. \href{https://bv.fapesp.br/en/bolsas/172826/effect-of-landscape-modifications-and-climate-changes-on-the-persistence-of-amphibians-in-the-atlant/}{\underline{Project link}}.
\item \textbf{Undergraduate research scholarship}, FAPESP, 2013-2015. Grant 2013/02883-7. \href{https://bv.fapesp.br/en/bolsas/142421/effect-of-fragmentation-on-the-persistence-of-anuran-amphibians-amphibia-anura-within-the-atlanti/}{\underline{Project link}}.
\end{itemize}
\end{rSection}

%------------------------------------------------------------------------
% Awards and Honors
%------------------------------------------------------------------------
\begin{rSection}{Awards and honors}
\begin{itemize}
\item \textbf{2026 SORTEE Commendation Award}, Society for Open, Reliable, and Transparent Ecology and Evolutionary Biology (SORTEE), 2026.
\item \textbf{The Harry R. Painton Award}, American Ornithological Society, 2023.
\item \textbf{High Academic Performance Award}, São Paulo State University (UNESP), 2015.
\end{itemize}
\end{rSection}

%-------------------------------------------------------------------------
% Publications
%-------------------------------------------------------------------------
\begin{rSection}{Publications}
\textbf{Selected peer-reviewed publications}
\begin{enumerate}
\item Gonçalves-Souza T, Chase J, Haddad N, \textbf{Vancine MH}, et al. 2025. \textbf{Species turnover does not rescue biodiversity in fragmented landscapes}. \textit{Nature} 640:702-706. \href{https://doi.org/10.1038/s41586-025-08688-7}{\underline{10.1038/s41586-025-08688-7}}
\item \textbf{Vancine MH}, Muylaert RL, Niebuhr BB, Oshima JEF, Tonetti V, Bernardo R, De Angelo C, Rosa MR, Grohmann CH, Ribeiro MC. 2024. \textbf{The Atlantic Forest of South America: spatiotemporal dynamics of remaining vegetation and implications for conservation}. \textit{Biological Conservation} 291:110499. \href{https://doi.org/10.1016/j.biocon.2024.110499}{\underline{10.1016/j.biocon.2024.110499}}
\item Marjakangas E, Abrego N, Grøtan V, et al., \textbf{Vancine MH}, et al. 2020. \textbf{Fragmented tropical forests lose mutualistic plant-animal interactions}. \textit{Diversity and Distributions} 26:154-168. \href{https://doi.org/10.1111/ddi.13010}{\underline{10.1111/ddi.13010}}
\item Santos JP, Sobral-Souza T, Brown KS Jr, \textbf{Vancine MH}, et al. 2020. \textbf{Effects of landscape modification on species richness patterns of fruit-feeding butterflies in Brazilian Atlantic Forest}. \textit{Diversity and Distributions} 26:196-208. \href{https://doi.org/10.1111/ddi.13007}{\underline{10.1111/ddi.13007}}
\item \textbf{Vancine MH}, Duarte KS, Souza YS, et al. 2018. \textbf{ATLANTIC AMPHIBIANS: a data set of amphibian communities from the Atlantic Forests of South America}. \textit{Ecology} 99:1692. \href{https://doi.org/10.1002/ecy.2392}{\underline{10.1002/ecy.2392}}
\item Sobral-Souza T, \textbf{Vancine MH}, Ribeiro MC, Lima-Ribeiro MS. 2018. \textbf{Efficiency of protected areas in Amazon and Atlantic Forest conservation: a spatio-temporal view}. \textit{Acta Oecologica} 87:1-7. \href{https://doi.org/10.1016/j.actao.2018.01.001}{\underline{10.1016/j.actao.2018.01.001}}
\item Beca G, \textbf{Vancine MH}, Carvalho CS, et al. 2017. \textbf{High mammal species turnover in forest patches immersed in biofuel plantations}. \textit{Biological Conservation} 210:352-359. \href{https://doi.org/10.1016/j.biocon.2017.02.033}{\underline{10.1016/j.biocon.2017.02.033}}
\item de Castro Pena JC, Goulart F, Fernandes GW, et al., \textbf{Vancine MH}, et al. 2017. \textbf{Impacts of mining activities on the potential geographic distribution of eastern Brazil mountaintop endemic species}. \textit{Perspectives in Ecology and Conservation} 15:172-178. \href{https://doi.org/10.1016/j.pecon.2017.07.005}{\underline{10.1016/j.pecon.2017.07.005}}
\end{enumerate}

\textbf{Books and book chapters}
\begin{itemize}
\item Da Silva FR, Gonçalves-Souza T, Paterno GB, Provete DB, \textbf{Vancine MH}. 2022. \textbf{Análises Ecológicas no R}. Nupeea / Canal 6. 640 p. \href{https://analises-ecologicas.com/}{\underline{link}}.
\item Lecours V, LeClerc E, Moudrý V, \textbf{Vancine MH}. 2026. \textbf{Ecological applications of geomorphometry}. In: \textit{Geomorphometry: Concepts, Software, Applications}. Elsevier. \href{https://doi.org/10.1016/B978-0-44-333376-7.00037-9}{\underline{DOI}}.
\end{itemize}

\textbf{Publication profiles}: \href{https://orcid.org/0000-0001-9650-7575}{\underline{ORCID}}, \href{https://www.webofscience.com/wos/author/record/837504}{\underline{Web of Science}}, \href{https://www.scopus.com/authid/detail.uri?authorId=57193451888}{\underline{Scopus}}, \href{https://scholar.google.com/citations?user=i-2xZBQAAAAJ}{\underline{Google Scholar}}.
\end{rSection}

%-------------------------------------------------------------------------
% Manuscripts under review
%-------------------------------------------------------------------------
\begin{rSection}{Manuscripts under review}
\begin{enumerate}
\item Morini RFS*, \textbf{Vancine MH*}, Martinez AN, Sales LP, Pires MM. \textbf{Habitat degradation threatens the topological diversity of seed-dispersal networks in the Atlantic Forest hotspot}. \textit{Global Change Biology}.
\item Redana M, Moura MR, Starinchak B, \textbf{Vancine MH}, Brando P, Goetz S, Jantz P. \textbf{Forest structure types in the southern Amazon Basin are influenced by anthropogenic edges and physiography}. \textit{Science of Remote Sensing}.
\item Martinez AN, \textbf{Vancine MH}, Silva W, Pizo MA, Keitt TH, Guimarães PR, Pires MM. \textbf{Estimating landscape connectivity for plants through the lens of seed dispersers}. \textit{Journal of Ecology}.
\item de Lima HVG, \textbf{Vancine MH}, Araújo HFP, Moura MR. \textbf{Identifying shortfalls in community ecological research within a high biodiversity and threatened hotspot}. \textit{Biological Conservation}.
\end{enumerate}
\end{rSection}

%-------------------------------------------------------------------------
% Teaching Experience
%-------------------------------------------------------------------------
\begin{rSection}{Teaching experience}
\textbf{Graduate courses}
\begin{itemize}
\item \textbf{Data analysis and graphical presentation using R (60 h)} \hfill{\em 2024, 2025, 2026}\\
Invited teacher, Ecology Graduate Program, State University of Campinas (Unicamp).
\item \textbf{Introduction to the use of geospatial data in R (60 h)} \hfill{\em 2020, 2021}\\
Invited teacher, Ecology, Biodiversity and Evolution Graduate Program, São Paulo State University (UNESP).
\item \textbf{Ecological Niche Modeling: theory and practice (45-60 h)} \hfill{\em 2016-2017}\\
Teaching assistant, graduate programs at UNESP and Unicamp.
\item \textbf{Introduction to Geoprocessing for Ethnobiology and Conservation (45 h)} \hfill{\em 2019}\\
Invited teacher, Rural Federal University of Pernambuco (UFRPE).
\end{itemize}

\textbf{Undergraduate teaching}
\begin{itemize}
\item \textbf{Statistical Models in Ecology (30 h)} \hfill{\em 2020}\\
Teacher, Ecology Undergraduate Program, UNESP.
\item \textbf{Quantitative Ecology (60 h)} \hfill{\em 2015}\\
Teaching assistant, Ecology Undergraduate Program, UNESP.
\end{itemize}

\textbf{Selected workshops}
\begin{itemize}
\item Introduction to R and ecological data analysis; species distribution modeling; geoprocessing with QGIS; version control with Git and GitHub; field herpetology. Workshops delivered at UNESP, UFRPE, Unicamp, and scientific meetings, 2014-2021.
\end{itemize}
\end{rSection}

%-------------------------------------------------------------------------
% Administrative experience
%-------------------------------------------------------------------------
\begin{rSection}{Administrative experience}
\begin{itemize}
\item \textbf{Student representative}, Graduate Program in Ecology, Evolution and Biodiversity, UNESP. \hfill{\em 2020-2021}
\end{itemize}
\end{rSection}

%-------------------------------------------------------------------------
% Professional affiliations and memberships
%-------------------------------------------------------------------------
\begin{rSection}{Professional affiliations and memberships}
\begin{itemize}
\item Associação Brasileira de Ecólogos (ABE) \hfill{\em 2018-present}
\item Associação Brasileira de Ciência Ecológica e Conservação (ABECO) \hfill{\em 2021-present}
\item Society for Open, Reliable, and Transparent Ecology and Evolutionary Biology (SORTEE) \hfill{\em 2025-present}
\end{itemize}
\end{rSection}

%-------------------------------------------------------------------------
% Supervision
%-------------------------------------------------------------------------
\begin{rSection}{Supervision}
\begin{itemize}
\item \textbf{Doctoral co-supervision}: Pâmela Rodrigues Braga (2026), Graduate Program in Ecology, Unicamp. Research on frugivory and seed dispersal by primates in the Amazon.
\item \textbf{Master's co-supervision}: Lucas de Souza Almeida (2022), Graduate Program in Zoology, UNESP. Research on the distribution and conservation of caecilians (Amphibia: Gymnophiona).
\item \textbf{Undergraduate research supervision/co-supervision}: Helena Thereza Carvalho de Oliveira (2021), Maria Eduarda Furlan (2025), and Kayane Ferreira Santana Trindade (2026), on amphibian reproductive patterns, rabies circulation and zoonotic risk, and sea-turtle habitat suitability.
\item \textbf{Undergraduate theses}: Helena Thereza Carvalho de Oliveira (2021), Bruno Eduardo Ribeiro Silva (2022), Maria Eduarda Furlan (2025), and Mariana Trindade Lahr (2025). Topics include amphibian ecology, species distribution modeling, zoonotic risk, and road ecology.
\end{itemize}
\end{rSection}

%-------------------------------------------------------------------------
% Services
%-------------------------------------------------------------------------
\begin{rSection}{Services}
\begin{itemize}
\item \textbf{Journal editor}: \textit{Austral Ecology}.
\item \textbf{Journal reviewer}: peer reviewer for journals in ecology, conservation, biogeography, zoology, and spatial analysis, including \textit{Nature}-family journals, \textit{Global Change Biology}, \textit{Journal of Biogeography}, \textit{Biological Conservation}, \textit{Landscape Ecology}, and \textit{Ecological Applications}.
\item \textbf{Funding and award reviewer}: São Paulo Research Foundation (FAPESP) and MapBiomas Award.
\end{itemize}
\end{rSection}

%-------------------------------------------------------------------------
% Skills
%-------------------------------------------------------------------------
\begin{rSection}{Skills}
\textbf{Programming and data analysis}
\begin{itemize}
\item R (advanced; tidyverse, Markdown, Quarto), Git/GitHub, LaTeX, Python, shell/bash; statistical analysis, spatial data handling and visualization, model selection, ordination, and machine learning.
\end{itemize}

\textbf{Geospatial software and platforms}
\begin{itemize}
\item QGIS and GRASS GIS (advanced); Google Earth Engine, ArcGIS, GeoDa, Fragstats, and GuidosToolbox.
\end{itemize}

\textbf{R packages and data resources}
\begin{itemize}
\item \textit{lsmetrics}: multiscale calculation and spatialization of landscape metrics using GRASS GIS. \href{https://mauriciovancine.github.io/lsmetrics}{\underline{Project page}}.
\item \textit{atlanticr}: access to datasets associated with the ATLANTIC data papers on the Atlantic Forest. \href{https://mauriciovancine.github.io/atlanticr}{\underline{Project page}}.
\item \textit{amphiBR}: data resource associated with the official Brazilian amphibian species list. \href{https://paulobarros.github.io/amphiBR}{\underline{Project page}}.
\item \textit{ecodados}: ecological database designed for teaching statistics. \href{https://paternogbc.github.io/ecodados/}{\underline{Project page}}.
\end{itemize}

\textbf{Languages}
\begin{itemize}
\item Portuguese (native); English (intermediate); Spanish (basic); Italian (basic).
\end{itemize}
\end{rSection}

%-------------------------------------------------------------------------
% Media
%-------------------------------------------------------------------------
\begin{rSection}{Media}
Selected media coverage of research on biodiversity, landscape fragmentation, and conservation:
\begin{itemize}
\item \textbf{Novo estudo destaca preocupação com estado de conservação da Mata Atlântica}. \textit{Agência FAPESP}, 2024.
\item \textbf{Desmatamento desacelera na Mata Atlântica, mas fragmentação aumenta}. \textit{Mongabay Brasil}, 2024.
\item \textbf{Protect one large forest, or many small ones? New study reignites conservation debate}. \textit{Mongabay}, 2025.
\item \textbf{Como uma base de dados ambientais espacializados ajuda na conservação da Mata Atlântica}. \textit{NEXO Acadêmico}, 2026.
\end{itemize}
\end{rSection}

%-------------------------------------------------------------------------
% References
%-------------------------------------------------------------------------
\begin{rSection}{References}
Available upon request. Full publication and citation records: \href{https://orcid.org/0000-0001-9650-7575}{\underline{ORCID}}, \href{https://scholar.google.com/citations?user=i-2xZBQAAAAJ}{\underline{Google Scholar}}.
\end{rSection}

\end{document}